% ============================================================
\chapter{Syst\`emes Non Lin\'eaires --- Lin\'earisation et Vari\'et\'e Centrale}
\label{ch:non_lineaires}
% ============================================================

Nous quittons maintenant le monde rassurant des syst\`emes lin\'eaires pour entrer dans celui, infiniment plus riche, des syst\`emes non lin\'eaires. Ici, les surprises abondent~: un pendule peut osciller ou tourner, une population peut exploser ou s'effondrer, un circuit \'electrique peut engendrer des oscillations auto-entretenues. Le th\'eor\`eme de Hartman-Grobman (1960) nous assure que, pr\`es d'un \'equilibre \emph{hyperbolique}, le portrait de phase non lin\'eaire ressemble qualitativement \`a celui du syst\`eme lin\'earis\'e. Mais que se passe-t-il lorsque des valeurs propres tombent sur l'axe imaginaire~? C'est l\`a qu'intervient la th\'eorie de la \emph{vari\'et\'e centrale}, d\'evelopp\'ee par Victor Pliss (1964) et Jack Carr (1981).

\section{Introduction}

Les syst\`emes non lin\'eaires pr\'esentent une richesse ph\'enom\'enologique
consid\'erable qui d\'epasse le cadre des syst\`emes lin\'eaires~:
multiplicit\'e de points d'\'equilibre, cycles limites, chaos. Pour
comprendre le comportement \emph{local} au voisinage d'un \'equilibre, on
recourt \`a la lin\'earisation (th\'eor\`eme de Hartman-Grobman) et, dans les
cas d\'eg\'en\'er\'es, \`a la th\'eorie des vari\'et\'es centrales.

\section{Le th\'eor\`eme de Hartman-Grobman}

\begin{theoreme}[Hartman-Grobman]
\label{thm:hartman-grobman}
Soit $\dot{x} = f(x)$ avec $f \in C^1$, $f(0) = 0$, et $A = Df(0)$.
Si $A$ est \textbf{hyperbolique} (aucune valeur propre sur l'axe imaginaire),
alors il existe un hom\'eomorphisme $h$ d\'efini dans un voisinage de $0$
tel que $h$ conjugue le flot de $\dot{x} = f(x)$ au flot lin\'eaire
$\dot{y} = Ay$ :
\[
    h(\varphi_t(x)) = e^{At}\, h(x)
\]
pour $t$ suffisamment petit et $x$ dans un voisinage de $0$.
\end{theoreme}

\begin{remarque}
Le th\'eor\`eme de Hartman-Grobman garantit que la topologie du portrait de
phase local est enti\`erement d\'etermin\'ee par la lin\'earisation lorsque
l'\'equilibre est hyperbolique. L'hom\'eomorphisme $h$ n'est en g\'en\'eral
pas diff\'erentiable.
\end{remarque}

\begin{definition}[\'Equilibre hyperbolique]
Un point d'\'equilibre $x^*$ est \textbf{hyperbolique} si la matrice
jacobienne $Df(x^*)$ n'a aucune valeur propre de partie r\'eelle nulle.
\end{definition}

\begin{definition}[\'Equilibre non hyperbolique]
Un \'equilibre est \textbf{non hyperbolique} (ou \textbf{d\'eg\'en\'er\'e})
si au moins une valeur propre de $Df(x^*)$ a une partie r\'eelle nulle.
Dans ce cas, la lin\'earisation ne d\'etermine pas la stabilit\'e : il faut
des outils suppl\'ementaires.
\end{definition}

\section{Vari\'et\'es invariantes locales}

\begin{theoreme}[Vari\'et\'es stable et instable]
\label{thm:varietes-stables}
Soit $x^* = 0$ un \'equilibre hyperbolique de $\dot{x} = f(x)$ avec
$A = Df(0)$. Soient $E^s$ et $E^u$ les sous-espaces stable et instable de
$A$. Alors il existe des vari\'et\'es locales :
\begin{itemize}
    \item $W^s_{\mathrm{loc}}(0)$ : \textbf{vari\'et\'e stable locale},
    tangente \`a $E^s$ en $0$, de dimension $\dim E^s$ ;
    \item $W^u_{\mathrm{loc}}(0)$ : \textbf{vari\'et\'e instable locale},
    tangente \`a $E^u$ en $0$, de dimension $\dim E^u$.
\end{itemize}
Ces vari\'et\'es sont de classe $C^k$ si $f \in C^k$, et sont localement
invariantes par le flot.
\end{theoreme}

\begin{definition}[Vari\'et\'es globales]
Les \textbf{vari\'et\'es stable et instable globales} sont d\'efinies par :
\[
    W^s(0) = \{x \in U : \varphi(t, x) \to 0 \text{ quand } t \to +\infty\},
\]
\[
    W^u(0) = \{x \in U : \varphi(t, x) \to 0 \text{ quand } t \to -\infty\}.
\]
\end{definition}

\begin{exemple}[Col en dimension 2]
Pour le syst\`eme $\dot{x} = -x + x^2$, $\dot{y} = y$, l'origine est un
col avec $E^s = \{y = 0\}$, $E^u = \{x = 0\}$. La vari\'et\'e stable locale
est la courbe $y = 0$, $\abs{x} < 1$, et la vari\'et\'e instable est l'axe $x = 0$.
\end{exemple}

\section{Th\'eorie de la vari\'et\'e centrale}

Lorsque l'\'equilibre est non hyperbolique, la dynamique essentielle est
captur\'ee par la \textbf{vari\'et\'e centrale}.

\begin{theoreme}[Existence de la vari\'et\'e centrale]
\label{thm:variete-centrale}
Soit $\dot{x} = f(x)$ avec $f(0) = 0$ et $A = Df(0)$. Supposons que $A$
ait $n_s$ valeurs propres \`a partie r\'eelle n\'egative, $n_u$ \`a partie
r\'eelle positive, et $n_c$ \`a partie r\'eelle nulle, avec
$n_s + n_u + n_c = n$. Alors il existe une vari\'et\'e locale
$W^c_{\mathrm{loc}}(0)$ :
\begin{itemize}
    \item de dimension $n_c$ ;
    \item tangente \`a $E^c$ en $0$ ;
    \item localement invariante par le flot ;
    \item de classe $C^k$ si $f \in C^k$ (mais pas n\'ecessairement analytique
    m\^eme si $f$ est analytique).
\end{itemize}
\end{theoreme}

\begin{attention}[title=\textbf{Non-unicit\'e de la vari\'et\'e centrale}]
Contrairement aux vari\'et\'es stable et instable, la vari\'et\'e centrale
n'est en g\'en\'eral \textbf{pas unique}. Deux vari\'et\'es centrales diff\'erentes
peuvent co\"incider \`a tout ordre mais diff\'erer par des termes
exponentiellement petits.
\end{attention}

\begin{theoreme}[R\'eduction \`a la vari\'et\'e centrale]
\label{thm:reduction-centrale}
Si $n_u = 0$ (pas de direction instable), la stabilit\'e de l'\'equilibre
$x^* = 0$ du syst\`eme complet est d\'etermin\'ee par la stabilit\'e de
$0$ dans la dynamique restreinte \`a la vari\'et\'e centrale $W^c$.
\end{theoreme}

\subsection{Calcul pratique de la vari\'et\'e centrale}

Consid\'erons le syst\`eme sous forme canonique :
\[
    \dot{x} = Cx + F(x, y), \qquad \dot{y} = Sy + G(x, y),
\]
o\`u $x \in \R^{n_c}$, $y \in \R^{n_s}$, $C$ a toutes ses valeurs propres
sur l'axe imaginaire, $S$ a toutes ses valeurs propres \`a partie r\'eelle
n\'egative, et $F$, $G$ sont d'ordre au moins 2.

La vari\'et\'e centrale est le graphe d'une fonction $y = h(x)$ avec
$h(0) = 0$, $Dh(0) = 0$. La fonction $h$ satisfait l'\'equation aux
d\'eriv\'ees partielles :
\[
    Dh(x)\,[Cx + F(x, h(x))] = Sh(x) + G(x, h(x)).
\]

En pratique, on d\'eveloppe $h$ en s\'erie de Taylor et on identifie les
coefficients ordre par ordre.

\begin{exemple}[Vari\'et\'e centrale en dimension 2]
Consid\'erons le syst\`eme
\[
    \dot{x} = xy, \qquad \dot{y} = -y + x^2.
\]
L'origine est un \'equilibre non hyperbolique : $A = \begin{pmatrix} 0 & 0 \\ 0 & -1 \end{pmatrix}$,
$E^c = \{y = 0\}$, $E^s = \{x = 0\}$.

On cherche $y = h(x) = ax^2 + bx^3 + \cdots$ Substituant dans l'\'equation
de la vari\'et\'e centrale :
\[
    h'(x) \cdot x\,h(x) = -h(x) + x^2.
\]
\`A l'ordre 2 : $0 = -ax^2 + x^2$, donc $a = 1$.
\`A l'ordre 3 : $2x \cdot x \cdot x^2 = -bx^3$, donc $2x^3$ versus
$-bx^3$ donne... recalculons : $h'(x) = 2x + 3bx^2 + \cdots$, et
$h'(x) \cdot x h(x) = (2x)(x)(x^2) + \cdots = 2x^4 + \cdots$ qui est
d'ordre 4. Donc $b = 0$ \`a l'ordre 3.

La dynamique sur la vari\'et\'e centrale est :
\[
    \dot{x} = x \cdot h(x) = x \cdot x^2 + \cdots = x^3 + \cdots
\]
Comme $\dot{x} = x^3 > 0$ pour $x > 0$ et $\dot{x} < 0$ pour $x < 0$,
l'\'equilibre est \textbf{instable}.
\end{exemple}

\section{Formes normales}

\begin{definition}[Forme normale de Poincar\'e]
Une \textbf{forme normale} est une simplification du syst\`eme non lin\'eaire
obtenue par des changements de variables polynomiaux successifs qui
\'eliminent les termes non r\'esonnants.
\end{definition}

\begin{theoreme}[Th\'eor\`eme de la forme normale]
Soit $\dot{x} = Ax + f_2(x) + f_3(x) + \cdots$ avec $f_k$ homog\`ene de
degr\'e $k$. Pour chaque $k \geq 2$, il existe un changement de variable
polynomial $x = y + \phi_k(y)$ qui \'elimine les termes non r\'esonnants
d'ordre $k$ dans $f_k$.
\end{theoreme}

\begin{definition}[R\'esonance]
Un mon\^ome $x^m = x_1^{m_1} \cdots x_n^{m_n}$ dans la $j$-\`eme composante
est \textbf{r\'esonnant} si
\[
    \langle m, \lambda \rangle - \lambda_j = 0,
\]
o\`u $\lambda = (\lambda_1, \ldots, \lambda_n)$ est le vecteur des valeurs
propres et $\langle m, \lambda \rangle = m_1\lambda_1 + \cdots + m_n\lambda_n$.
\end{definition}

\section{Lin\'earisation diff\'erentiable}

\begin{theoreme}[Sternberg]
Si $x^* = 0$ est hyperbolique et si les valeurs propres de $Df(0)$ ne
satisfont aucune relation de r\'esonance, alors il existe un
diff\'eomorphisme local de classe $C^\infty$ qui conjugue le syst\`eme
non lin\'eaire au syst\`eme lin\'eaire $\dot{y} = Ay$.
\end{theoreme}

\begin{remarque}
Le th\'eor\`eme de Sternberg am\'eliore celui de Hartman-Grobman en donnant
un diff\'eomorphisme (et pas seulement un hom\'eomorphisme) sous des
conditions de non-r\'esonance.
\end{remarque}

\section{Application : dynamique de population}

\begin{exemple}[Comp\'etition entre deux esp\`eces]
Le mod\`ele de comp\'etition de Lotka-Volterra est
\[
    \dot{x}_1 = x_1(1 - x_1 - \alpha x_2), \qquad
    \dot{x}_2 = x_2(1 - x_2 - \beta x_1),
\]
avec $\alpha, \beta > 0$. Les points d'\'equilibre sont $(0,0)$, $(1,0)$,
$(0,1)$ et, si $\alpha \neq 1$ et $\beta \neq 1$, le point de coexistence
\[
    x^* = \left(\frac{1-\alpha}{1-\alpha\beta}, \frac{1-\beta}{1-\alpha\beta}\right).
\]
La lin\'earisation en chaque point d'\'equilibre d\'etermine la dynamique locale.
\end{exemple}

\begin{codeblock}[title=\textbf{Vari\'et\'es stable et instable du col}]
\begin{minted}{python}
import numpy as np
from scipy.integrate import solve_ivp
import matplotlib.pyplot as plt

def system(t, state):
    x, y = state
    return [x - x**2, -y + x*y]

fig, ax = plt.subplots(figsize=(8, 6))

# Stable manifold (integrate backward from near saddle)
for eps in [0.01, -0.01]:
    sol = solve_ivp(lambda t, s: [-v for v in system(t, s)],
                    [0, 5], [0, eps], max_step=0.01)
    ax.plot(sol.y[0], sol.y[1], 'b-', lw=2, label='$W^s$' if eps > 0 else '')

# Unstable manifold (integrate forward)
for eps in [0.01, -0.01]:
    sol = solve_ivp(system, [0, 5], [eps, 0], max_step=0.01)
    ax.plot(sol.y[0], sol.y[1], 'r-', lw=2, label='$W^u$' if eps > 0 else '')

ax.set_xlim(-1, 2); ax.set_ylim(-1, 2)
ax.set_xlabel('$x$'); ax.set_ylabel('$y$')
ax.legend(); ax.grid(True, alpha=0.3)
ax.set_title("Varietes stable et instable")
plt.tight_layout()
plt.savefig("stable_unstable_manifolds.pdf")
plt.show()
\end{minted}
\end{codeblock}

\section{Exercices}

\begin{exercice}[Hartman-Grobman]
V\'erifier que l'\'equilibre $(0,0)$ du syst\`eme
$\dot{x} = -x + y^2$, $\dot{y} = -2y + x^2$ est hyperbolique.
Quel est le type du portrait de phase local ?
\end{exercice}

\begin{exercice}[Vari\'et\'e centrale]
Calculer la vari\'et\'e centrale et d\'eterminer la stabilit\'e de l'origine pour
\[
    \dot{x} = -x^2y, \qquad \dot{y} = -y + x^2.
\]
\end{exercice}

\begin{exercice}[Bifurcation transcritique]
\'Etudier le syst\`eme $\dot{x} = \mu x - x^2$, $\dot{y} = -y$ au voisinage
de $(\mu, x, y) = (0, 0, 0)$ par r\'eduction \`a la vari\'et\'e centrale.
\end{exercice}

\begin{exercice}[Forme normale]
Mettre sous forme normale \`a l'ordre 3 le syst\`eme
$\dot{x} = -y + ax^2 + bxy$, $\dot{y} = x + cx^2 + dxy$.
\end{exercice}

\begin{exercice}[Vari\'et\'es globales]
Tracer num\'eriquement les vari\'et\'es stable et instable de l'\'equilibre
$(1, 0)$ du syst\`eme $\dot{x} = x(1-x) - xy$, $\dot{y} = y(x - 1)$.
\end{exercice}

\begin{exercice}[Non-unicit\'e]
Montrer que pour le syst\`eme $\dot{x} = x^2$, $\dot{y} = -y$, la vari\'et\'e
centrale n'est pas unique en exhibant deux fonctions $h_1(x)$ et $h_2(x)$
satisfaisant l'\'equation de la vari\'et\'e centrale.
\end{exercice}
